%!TeX root=../persuasiontop.tex
\chapter[Chapter \thechapter]{}

\begin{figure}[t!]
\centering
\includegraphics[width=\linewidth]{nouveau02}
\end{figure}

\lettrine[lraise=0.3]{M}{r} Shepherd, a civil, cautious lawyer, who, whatever might be his hold or his views on Sir Walter, would rather have the \textit{disagreeable} prompted by anybody else, excused himself from offering the slightest hint, and only begged leave to recommend an implicit reference to the excellent judgement of Lady Russell, from whose known good sense he fully expected to have just such resolute measures advised as he meant to see finally adopted.

Lady Russell was most anxiously zealous on the subject, and gave it much serious consideration. She was a woman rather of sound than of quick abilities, whose difficulties in coming to any decision in this instance were great, from the opposition of two leading principles. She was of strict integrity herself, with a delicate sense of honour; but she was as desirous of saving Sir Walter's feelings, as solicitous for the credit of the family, as aristocratic in her ideas of what was due to them, as anybody of sense and honesty could well be. She was a benevolent, charitable, good woman, and capable of strong attachments, most correct in her conduct, strict in her notions of decorum, and with manners that were held a standard of good-breeding. She had a cultivated mind, and was, generally speaking, rational and consistent—but she had prejudices on the side of ancestry; she had a value for rank and consequence, which blinded her a little to the faults of those who possessed them. Herself the widow of only a knight, she gave the dignity of a baronet all its due; and Sir Walter, independent of his claims as an old acquaintance, an attentive neighbour, an obliging landlord, the husband of her very dear friend, the father of Anne and her sisters, was, as being Sir Walter, in her apprehension, entitled to a great deal of compassion and consideration under his present difficulties.

They must retrench; that did not admit of a doubt. But she was very anxious to have it done with the least possible pain to him and Elizabeth. She drew up plans of economy, she made exact calculations, and she did what nobody else thought of doing: she consulted Anne, who never seemed considered by the others as having any interest in the question. She consulted, and in a degree was influenced by her in marking out the scheme of retrenchment which was at last submitted to Sir Walter. Every emendation of Anne's had been on the side of honesty against importance. She wanted more vigorous measures, a more complete reformation, a quicker release from debt, a much higher tone of indifference for everything but justice and equity.

<If we can persuade your father to all this,> said Lady Russell, looking over her paper, <much may be done. If he will adopt these regulations, in seven years he will be clear; and I hope we may be able to convince him and Elizabeth, that Kellynch Hall has a respectability in itself which cannot be affected by these reductions; and that the true dignity of Sir Walter Elliot will be very far from lessened in the eyes of sensible people, by acting like a man of principle. What will he be doing, in fact, but what very many of our first families have done, or ought to do? There will be nothing singular in his case; and it is singularity which often makes the worst part of our suffering, as it always does of our conduct. I have great hope of prevailing. We must be serious and decided; for after all, the person who has contracted debts must pay them; and though a great deal is due to the feelings of the gentleman, and the head of a house, like your father, there is still more due to the character of an honest man.>

This was the principle on which Anne wanted her father to be proceeding, his friends to be urging him. She considered it as an act of indispensable duty to clear away the claims of creditors with all the expedition which the most comprehensive retrenchments could secure, and saw no dignity in anything short of it. She wanted it to be prescribed, and felt as a duty. She rated Lady Russell's influence highly; and as to the severe degree of self-denial which her own conscience prompted, she believed there might be little more difficulty in persuading them to a complete, than to half a reformation. Her knowledge of her father and Elizabeth inclined her to think that the sacrifice of one pair of horses would be hardly less painful than of both, and so on, through the whole list of Lady Russell's too gentle reductions.

How Anne's more rigid requisitions might have been taken is of little consequence. Lady Russell's had no success at all: could not be put up with, were not to be borne. <What! every comfort of life knocked off! Journeys, London, servants, horses, table—contractions and restrictions every where! To live no longer with the decencies even of a private gentleman! No, he would sooner quit Kellynch Hall at once, than remain in it on such disgraceful terms.>

<Quit Kellynch Hall.> The hint was immediately taken up by Mr Shepherd, whose interest was involved in the reality of Sir Walter's retrenching, and who was perfectly persuaded that nothing would be done without a change of abode. <Since the idea had been started in the very quarter which ought to dictate, he had no scruple,> he said, <in confessing his judgement to be entirely on that side. It did not appear to him that Sir Walter could materially alter his style of living in a house which had such a character of hospitality and ancient dignity to support. In any other place Sir Walter might judge for himself; and would be looked up to, as regulating the modes of life in whatever way he might choose to model his household.>

Sir Walter would quit Kellynch Hall; and after a very few days more of doubt and indecision, the great question of whither he should go was settled, and the first outline of this important change made out.

There had been three alternatives, London, Bath, or another house in the country. All Anne's wishes had been for the latter. A small house in their own neighbourhood, where they might still have Lady Russell's society, still be near Mary, and still have the pleasure of sometimes seeing the lawns and groves of Kellynch, was the object of her ambition. But the usual fate of Anne attended her, in having something very opposite from her inclination fixed on. She disliked Bath, and did not think it agreed with her; and Bath was to be her home.

Sir Walter had at first thought more of London; but Mr Shepherd felt that he could not be trusted in London, and had been skilful enough to dissuade him from it, and make Bath preferred. It was a much safer place for a gentleman in his predicament: he might there be important at comparatively little expense. Two material advantages of Bath over London had of course been given all their weight: its more convenient distance from Kellynch, only fifty miles, and Lady Russell's spending some part of every winter there; and to the very great satisfaction of Lady Russell, whose first views on the projected change had been for Bath, Sir Walter and Elizabeth were induced to believe that they should lose neither consequence nor enjoyment by settling there.

Lady Russell felt obliged to oppose her dear Anne's known wishes. It would be too much to expect Sir Walter to descend into a small house in his own neighbourhood. Anne herself would have found the mortifications of it more than she foresaw, and to Sir Walter's feelings they must have been dreadful. And with regard to Anne's dislike of Bath, she considered it as a prejudice and mistake arising, first, from the circumstance of her having been three years at school there, after her mother's death; and secondly, from her happening to be not in perfectly good spirits the only winter which she had afterwards spent there with herself.

Lady Russell was fond of Bath, in short, and disposed to think it must suit them all; and as to her young friend's health, by passing all the warm months with her at Kellynch Lodge, every danger would be avoided; and it was in fact, a change which must do both health and spirits good. Anne had been too little from home, too little seen. Her spirits were not high. A larger society would improve them. She wanted her to be more known.

The undesirableness of any other house in the same neighbourhood for Sir Walter was certainly much strengthened by one part, and a very material part of the scheme, which had been happily engrafted on the beginning. He was not only to quit his home, but to see it in the hands of others; a trial of fortitude, which stronger heads than Sir Walter's have found too much. Kellynch Hall was to be let. This, however, was a profound secret, not to be breathed beyond their own circle.

Sir Walter could not have borne the degradation of being known to design letting his house. Mr Shepherd had once mentioned the word <advertise,> but never dared approach it again. Sir Walter spurned the idea of its being offered in any manner; forbad the slightest hint being dropped of his having such an intention; and it was only on the supposition of his being spontaneously solicited by some most unexceptionable applicant, on his own terms, and as a great favour, that he would let it at all.

How quick come the reasons for approving what we like! Lady Russell had another excellent one at hand, for being extremely glad that Sir Walter and his family were to remove from the country. Elizabeth had been lately forming an intimacy, which she wished to see interrupted. It was with the daughter of Mr Shepherd, who had returned, after an unprosperous marriage, to her father's house, with the additional burden of two children. She was a clever young woman, who understood the art of pleasing—the art of pleasing, at least, at Kellynch Hall; and who had made herself so acceptable to Miss Elliot, as to have been already staying there more than once, in spite of all that Lady Russell, who thought it a friendship quite out of place, could hint of caution and reserve.

Lady Russell, indeed, had scarcely any influence with Elizabeth, and seemed to love her, rather because she would love her, than because Elizabeth deserved it. She had never received from her more than outward attention, nothing beyond the observances of complaisance; had never succeeded in any point which she wanted to carry, against previous inclination. She had been repeatedly very earnest in trying to get Anne included in the visit to London, sensibly open to all the injustice and all the discredit of the selfish arrangements which shut her out, and on many lesser occasions had endeavoured to give Elizabeth the advantage of her own better judgement and experience; but always in vain: Elizabeth would go her own way; and never had she pursued it in more decided opposition to Lady Russell than in this selection of Mrs Clay; turning from the society of so deserving a sister, to bestow her affection and confidence on one who ought to have been nothing to her but the object of distant civility.

From situation, Mrs Clay was, in Lady Russell's estimate, a very unequal, and in her character she believed a very dangerous companion; and a removal that would leave Mrs Clay behind, and bring a choice of more suitable intimates within Miss Elliot's reach, was therefore an object of first-rate importance.